\section{总结与延伸}

\subsection{全课知识图谱}

本课建立了一条从数值表示到训练策略的完整认知链：

\begin{center}
\begin{tikzpicture}[node distance=0.9cm, every node/.style={draw, rounded corners, minimum width=4cm, minimum height=0.8cm, font=\small}]
\node (fp) {浮点数表示};
\node (mp) [below=of fp] {混合精度训练};
\node (ddp) [below=of mp] {DDP 数据并行};
\node (zero) [below=of ddp] {ZeRO / FSDP};
\node (lora) [below=of zero] {LoRA 参数高效微调};
\draw[->, thick] (fp) -- (mp);
\draw[->, thick] (mp) -- (ddp);
\draw[->, thick] (ddp) -- (zero);
\draw[->, thick] (zero) -- (lora);
\node[draw=none, right=1cm of fp, text width=5.5cm, font=\scriptsize] {FP32/FP16/BF16 的指数位与尾数位权衡};
\node[draw=none, right=1cm of mp, text width=5.5cm, font=\scriptsize] {Master Weights + Loss Scaling / BF16 AutoCast};
\node[draw=none, right=1cm of ddp, text width=5.5cm, font=\scriptsize] {All-Reduce 同步梯度，内存无优化};
\node[draw=none, right=1cm of zero, text width=5.5cm, font=\scriptsize] {分片优化器/梯度/参数，Reduce-Scatter + All-Gather};
\node[draw=none, right=1cm of lora, text width=5.5cm, font=\scriptsize] {低秩分解 $\Delta W = BA$，冻结原参数};
\end{tikzpicture}
\end{center}

\subsection{关键 Takeaways}

\begin{importantbox}{五条核心原则}
\begin{enumerate}
    \item \textbf{始终启用混合精度训练}：BF16（Ampere+）或 FP16 + GradScaler，几乎没有质量损失
    \item \textbf{优化器状态是显存大户}：Adam 的 momentum 和 variance 占 12 字节/参数，远超模型本身
    \item \textbf{ZeRO Stage 1/2 免费}：利用 All-Reduce = Reduce-Scatter + All-Gather 的等价性，不增加通信即可节省显存
    \item \textbf{FSDP 支持超大模型}：有额外通信开销，但可通过计算-通信重叠来隐藏
    \item \textbf{LoRA 是最后的利器}：当全量微调不可行时，低秩适配以极少的参数达到接近全量微调的效果
\end{enumerate}
\end{importantbox}

\subsection{拓展阅读}

\begin{itemize}
    \item Micikevicius et al., 2018. \textit{Mixed Precision Training}. ICLR 2018. \url{https://arxiv.org/abs/1710.03740}
    \item Rajbhandari et al., 2020. \textit{ZeRO: Memory Optimizations Toward Training Trillion Parameter Models}. SC 2020. \url{https://arxiv.org/abs/1910.02054}
    \item Hu et al., 2021. \textit{LoRA: Low-Rank Adaptation of Large Language Models}. ICLR 2022. \url{https://arxiv.org/abs/2106.09685}
    \item PyTorch FSDP 文档：\url{https://pytorch.org/docs/stable/fsdp.html}
    \item PyTorch AMP 文档：\url{https://pytorch.org/docs/stable/amp.html}
    \item Hugging Face PEFT 库：\url{https://github.com/huggingface/peft}
    \item DeepSpeed ZeRO 文档：\url{https://www.deepspeed.ai/tutorials/zero/}
\end{itemize}

\end{document}
